\documentclass{article}
\usepackage[T1]{fontenc}
\usepackage{lmodern}
\usepackage{graphicx}
\usepackage{amsmath,amssymb}
\usepackage[a4paper,margin=2cm]{geometry}
\usepackage[absolute,overlay]{textpos}
\setlength{\TPHorizModule}{1mm}
\setlength{\TPVertModule}{1mm}
\usepackage[indonesian]{babel}
\usepackage{hyperref}
\setlength{\emergencystretch}{2em}
% unit: Halaman 48

\begin{document}

% === Halaman 48 (python/native) ===

• Metode kedua adalah dengan metode Kolbitz. Langkah-langkahnya adalah 

sebagai berikut:

   1.  Pilih sebuah kurva eliptik y2  x3 + ax + b  (mod p) yang mengandung N buah 

titik.

   2.  Misalkan karakter-karakter penyusun pesan adalah angka 0, 1, 2, …, 9 dan 

huruf A, B, C, … Z yang dikodekan menjadi 10, 11, …, 35.

   3.  Kodekan setiap karakater di dalam pesan menjadi nilai m di antara 0 dan 35.

   4.  Pilih sebuah bilangan bulat k sebagai parameter basis (disepakati kedua pihak). 

   5.  Untuk setiap nilai mk, nyatakan x = mk + 1, sulihkan x ke dalam y2  x3 + ax + b  

(mod p) lalu tentukan nilai y yang memenuhi.

   6.  Jika tidak ada nilai y yang memenuhi, coba untuk x = mk + 2, x = mk + 3, dst, 

sampai y2  x3 + ax + b  (mod p) dapat dipecahkan.

   7 . Pada proses decoding, untuk titik (x, y), tentukan nilai m terbesar tetapi lebih 

kecil dari (x – 1)/k. Kodekan titik (x, y) menjadi symbol m.  

48

\end{document}
